\section{Analytical methods and mathematical scope}

The comparison uses one recorded customer cohort, feature contract, set of candidate grids, deployment constraints and evaluation measures. K Means provides a centre-based reference partition that can be approximated by ExKMC rules. CLASSIX provides a separate clustering result with native aggregation provenance. ExKMC is used here to approximate a fitted K Means reference, not as an independently selected customer clustering method; no composite score ranks it against CLASSIX. The same change in represented information is traced through each partition and its corresponding explanation.

\subsection{Clustering candidates and selection}

\paragraph{K Means: a partition around centres.}
For $n$ observations $x_i\in\mathbb{R}^d$ ($i=1,\ldots,n$), $d$ denotes the number of features. K Means seeks assignments $z_i\in\{1,\ldots,K\}$ and centres $\mu_j\in\mathbb{R}^d$ ($j=1,\ldots,K$) that minimise
\[
 \min_{z_1,\ldots,z_n,\,\mu_1,\ldots,\mu_K}
 \sum_{i=1}^{n}\lVert x_i-\mu_{z_i}\rVert_2^2.
\]
The Euclidean norm $\lVert\cdot\rVert_2$ measures straight-line distance; squaring it gives more weight to large deviations. The usual alternating procedure assigns each observation to its nearest centre, then replaces each centre by the mean of its assigned observations \citep{macqueen1967kmeans,sklearn172}. Each step does not increase the objective, but the fitted solution need not be the global minimum. The nearest-centre regions are Voronoi cells: each contains the points closer to its centre than to any other. k-means++ initialisation favours points far from already chosen centres using squared-distance sampling; repeated initialisations reduce dependence on one starting configuration \citep{arthur2007kmeanspp}.

Feature scaling matters because every coordinate contributes to squared distance. Adding nine standardised behavioural dimensions changes that sum even when the original three coordinates are unchanged. Centres and their nearest-centre assignments can consequently change for the same purchasers. Standardisation removes measurement-unit differences; it does not make the choice of features inconsequential.

\paragraph{CLASSIX: local groups followed by merging.}
CLASSIX first aggregates observations into small groups and then joins groups into clusters \citep{chen2024classix}. The study supplies standardised observations. In the version 1.5.1 Euclidean PCA branch, the package additionally centres them and divides by the median centred Euclidean norm (using one if that scale is zero). The radius below is in these processed coordinates, not in raw feature units \citep{classix151}.

The \emph{principal sorting direction} is a unit vector along which the processed observations have maximum variance. An observation's \emph{projected score} is its scalar coordinate along that direction. After sorting by these scores, the first unassigned observation becomes a \emph{starting point}. Later unassigned observations within the \emph{aggregation radius} of that point join its \emph{aggregation group}. Points already assigned to an earlier group are skipped. The next unassigned point starts another group, until every observation belongs to one. A starting point is therefore an observed representative, not a fitted centroid; a group is an intermediate object, not necessarily a final cluster.

For \emph{distance merging}, groups are linked when their starting points are within the radius multiplied by the merge scale, and connected groups share a cluster. For \emph{density merging}, the criterion uses counts per volume in overlapping radius balls: version 1.5.1 links a pair when the intersection density is at least one of the two ball densities. It thus tests local overlap density rather than only centre separation. Minimum cluster size controls the subsequent treatment of small merged components. The recorded distance branch reassigns groups from undersized components to a nearest eligible component. In the density branch, disabling post allocation leaves rejected observations labelled $-1$, or \emph{noise}. Noise denotes absence from a retained cluster under that configuration, not an invalid customer record. Small-group merging and final outlier allocation are distinct options \citep{classix151}.

The resulting \emph{provenance} connects each observation to its aggregation group, starting point and final label. This construction record is the reason for using CLASSIX here. It explains its own grouping process, whereas a post hoc tree approximates an already fitted reference. Neither representation is preferred in advance.

K Means uses k-means++ seeding, twenty initialisations and the fixed seed 20260801 \citep{macqueen1967kmeans,arthur2007kmeanspp}. The number of clusters ranges from 2 to 12. CLASSIX uses Euclidean distance and sorting along the first principal direction. Radius ranges from 0.25 to 1.75 in increments of 0.05. Minimum cluster size takes values 1, 5 and 10, while group merging uses either distance or density. The merge scale is 1.5, tiny groups can be merged, and post allocation is disabled. Each representation has 11 K Means candidates and 186 CLASSIX candidates. A recorded boundary diagnostic triggered a symmetric radius extension from 1.80 to 2.25 for both representations. The extension tests grid sensitivity and does not replace the primary selection.

Each method and representation reports an unconstrained silhouette optimum and a deployment-acceptable result. Silhouette compares an observation's mean distance to its own cluster with its smallest mean distance to another cluster; larger values indicate better separation relative to cohesion. Davies--Bouldin averages each cluster's worst ratio of within-cluster spread to centre separation, so smaller values are preferred. Calinski--Harabasz compares between-cluster with within-cluster squared dispersion, adjusted for sample and cluster counts, and larger values are preferred. These are different geometric criteria, not competing estimates of ground truth. A candidate is internally valid only when at least two nonnoise clusters remain and silhouette, Davies Bouldin and Calinski Harabasz indices are finite \citep{rousseeuw1987silhouettes,davies1979separation,calinski1974dendrite}. Deployment acceptability further requires 2 to 12 nonnoise clusters, no more than 20 per cent noise and no single nonnoise cluster larger than 95 per cent of nonnoise observations. Selection maximises silhouette, then prefers lower Davies Bouldin, higher Calinski Harabasz, a smaller largest cluster share, less noise, fewer clusters and finally the lexical parameter record. Runtime is measured but never selects a model. Keeping the unconstrained result reveals when an attractive internal score is produced by a partition that is unsuitable under the stated operational constraints.

Internal indices exclude observations labelled as noise, so CLASSIX index comparisons refer to each partition's nonnoise subset. Cross-representation partition comparisons retain every purchaser because a noise assignment is itself a difference in assignment. Fixed K comparisons use the same value from 2 to 12. Fixed CLASSIX comparisons use the same radius, minimum size and merge rule. Independently selected deployment partitions are compared by adjusted Rand index, arithmetic-mean normalised mutual information and variation of information \citep{hubert1985partitions,meila2007vi}. ARI measures agreement in whether pairs of observations share a cluster, adjusted for the agreement expected under its fixed-size random-partition model. One means identical partitions up to labels; zero is that chance-adjusted reference, and negative values are possible. It can compare two fitted partitions without ground truth: here it measures agreement, not accuracy. Arithmetic-mean NMI divides shared label information by the mean entropy of the two partitions; VI is their sum of conditional entropies, measured in nats with natural logarithms. Higher NMI and lower VI indicate closer agreement \citep{sklearn172,meila2007vi}. A Hungarian match finds the one-to-one relabelling that maximises overlapping assignments for transition tables \citep{kuhn1955hungarian}. The resulting migration share describes changed assignments and does not imply that numeric labels identify persistent customer types.

\subsection{Resampling and explanation evaluation}

Fixed-parameter overlap-resampling stability uses ten pairs of independent 80 per cent samples without replacement. The same seed pairs are used for both representations and methods. Each sample refits its transformations and clustering with the full-data deployment parameters. Adjusted Rand index is calculated only on purchaser identifiers present in both members of a pair. The statistic measures sensitivity to sampled customers and, for K Means, initialisation, conditional on the selected configuration. It does not measure stability of model selection, calendar time or an external cohort.

The CLASSIX explanation retains public aggregation group assignments, starting point indices, final labels and cluster sizes. Merge membership is reconstructed from public group and final assignments because version 1.5.1 does not expose a documented complete edge list. Native aggregation groups use sorted rows, so the export restores input order with the saved inverse permutation before linking groups to visitors. Every exported provenance label is checked against the selected main result for the same visitor, and every representative must belong to its reported group. The trace is an explanation of the fitted CLASSIX partition and does not claim that a group boundary is a human validated customer rule.

ExKMC approximates K Means labels under five deterministic 80 per cent training and 20 per cent test splits. The primary full cohort selection gives four K Means clusters in both representations, so K is fixed at four before splitting. Log transformations, standardisation, the K Means reference and the ExKMC tree are fitted using training rows only. The test target is the label returned by the fitted K Means predictor. This estimates explanation fidelity conditional on the primary K choice rather than nested selection of K. A leaf is a terminal decision-tree region; the leaf budget caps the number of such regions, and tree depth counts tests on the longest root-to-leaf path. The leaf budget ranges from K to four times K. The main comparison fixes the budget at K, while the wider range measures the fidelity and complexity tradeoff. A supplementary analysis selects K from 2 to 12 using each training split alone. Fidelity is the proportion of exact label agreements. Complexity is reported through effective leaves, depth, path conditions and features. Rule thresholds are exported in both standardised and original units.

Technical addressability means satisfying the study's specified size, coverage and named-feature distinctiveness requirements. It is evaluated for selected CLASSIX and K Means clusters through a documented audit rather than subjective cluster names. A nonnoise cluster must contain at least 50 customers and one per cent of nonnoise customers, retain at least 80 per cent category coverage, and differ from the global median by at least half an interquartile range on two to four named features. Passing this audit shows that a cluster can be expressed as a limited set of query conditions with adequate coverage. It does not demonstrate intervention effectiveness, revenue change or human comprehension.

\subsection{Projection filtering and complexity}

CLASSIX can rule out some candidate neighbours before computing their full Euclidean distance. The purpose of this section is to justify that computational shortcut and distinguish its theoretical cost from recorded performance. Proposition 3.1 restates the projection argument underlying CLASSIX aggregation \citep{chen2024classix}; it is not a new clustering objective.

Let $x_i\in\mathbb{R}^d$ be processed observation $i$, where $d$ is the number of coordinates, and let $u$ be a unit sorting direction. Define its score by $s_i=u^{\mathsf T}x_i$. A starting point is denoted by $c$, with score $s(c)=u^{\mathsf T}c$, and $r>0$ is the aggregation radius in the same processed coordinates. Scores are traversed in nondecreasing order.

\begin{proposition}
\label{prop:projection}
If a later point satisfies $s_i-s(c)>r$, then it cannot satisfy $\lVert x_i-c\rVert_2\leq r$.
\end{proposition}
\begin{proof}
The Cauchy--Schwarz inequality bounds an inner product by the product of the two vector lengths. Since the projection direction has length one,
\[
 |s_i-s(c)|=|u^{\mathsf T}(x_i-c)|
 \leq \lVert u\rVert_2\lVert x_i-c\rVert_2
 =\lVert x_i-c\rVert_2.
\]
Thus a projected separation greater than $r$ implies a full distance greater than $r$.
\end{proof}

In words, points already too far apart along one unit direction cannot be close in the full space. Once a sorted score exceeds the window, every later score does too, so the scan for this starting point can stop. The converse fails: a point inside the score window may be far away in other coordinates and still needs an exact distance test. Equality remains eligible in the implementation. The argument is exact in real arithmetic; it does not remove floating-point boundary effects or the order dependence of greedy aggregation.

Let $n$ be the number of observations and $q$ the number of resulting aggregation groups. At most $n(n-1)/2$ distinct later-point pairs can be tested during aggregation. Distance merging has at most $q(q-1)/2$ distinct starting-point pairs. Each full distance involves $d$ coordinates, giving the conservative distance-arithmetic bound
\[
 O\bigl((n^2+q^2)d\bigr).
\]
This is a worst-case bound for those distance operations, not an average runtime law. Principal-direction estimation, scalar sorting and bookkeeping have additional costs; small-component reassignment can add a quadratic group-distance pass. The bound on distinct pairs need not equal the number of vectorised arithmetic operations, which may include self-comparisons.

The saved package counter counts retained aggregation distance comparisons on the tested inputs; elapsed runtime also includes preprocessing, merging and implementation overhead. Neither an observed small counter nor a short runtime proves linear complexity. The distinction matters here because the grouping trace arises during clustering itself. Its computational mechanism differs from the later construction of an ExKMC tree.

\subsection{A local stability certificate}
\label{sec:local_certificate}

Section~\ref{sec:perturbation_results} changes numerical feature values slightly and measures whether CLASSIX assignments change. One possible cause of change is rotation of the principal sorting direction. This section asks a narrower mathematical question: under what sufficient conditions do the finite decisions of a distance-merging fit remain unchanged? It separates stability of a direction from stability of a complete partition.

Write $S$ for the covariance matrix of the centred, fully processed observations, and $\lambda_1\geq\lambda_2$ for its two largest eigenvalues. The \emph{eigengap} $\gamma=\lambda_1-\lambda_2>0$ measures how clearly the leading direction is separated from the next. A perturbation is a change of the numerical input: it gives processed observations $\widehat x_i$ and covariance $S+\Delta S$, where $\Delta S$ is the covariance change. Let $u$ and $\widehat u$ be unit leading eigenvectors before and after this change. Choose their relative signs so that $u^{\mathsf T}\widehat u\geq0$.

Davis--Kahan bounds the rotation of an eigenspace when its matrix changes \citep{davis1970rotation}. Using the one-dimensional form of \citet[Corollary 1]{yu2015variant}, and restricting attention to $2\lVert\Delta S\rVert_2<\gamma$, gives
\[
 \sin\angle(\widehat u,u)\leq
 \frac{2\lVert\Delta S\rVert_2}{\gamma}.
\]
Here the matrix norm is the spectral norm (the largest singular value); the angle is the acute angle between the two directions. A larger eigengap permits a smaller rotation bound for the same covariance change. If the leading directions are nearly tied, this argument gives little protection.

Let $\varepsilon_x$ bound each processed displacement $\lVert\widehat x_i-x_i\rVert_2$, and let $B$ bound every original processed norm $\lVert x_i\rVert_2$. The sign-aligned unit vectors obey $\lVert\widehat u-u\rVert_2\leq\sqrt{2}\sin\angle(\widehat u,u)$. Consequently,
\begin{align*}
 |\widehat u^{\mathsf T}\widehat x_i-u^{\mathsf T}x_i|
 &\leq \lVert\widehat x_i-x_i\rVert_2
       +\lVert\widehat u-u\rVert_2\lVert x_i\rVert_2\\
 &\leq \varepsilon_x+\sqrt{2}B\frac{2\lVert\Delta S\rVert_2}{\gamma}
 =:b_s.
\end{align*}
Thus $b_s$ is a maximum score movement, accounting for movement of both the point and the direction. The input to this bound is after all preprocessing: changes of fitted centring and scale must be included in $\varepsilon_x$ and $\Delta S$. A Gaussian noise standard deviation is not itself a bound on every processed displacement.

The remaining conditions protect individual decisions from crossing thresholds. Define $g_s$ as the minimum adjacent gap in the original sorted scores. Define $g_o=|u^{\mathsf T}x_1|$ for the first input row used by version 1.5.1's score-sign convention. Define $g_w$ as the smallest absolute distance of any relevant score difference from its active search-window threshold, including both aggregation radius $r$ and distance-merge radius $1.5r$. Let $g_r$ be the smallest absolute margin of an aggregation distance from $r$, and $g_m$ the corresponding distance-merge margin from $1.5r$. These minima refer to the decisions traversed by the original fit; they must be positive for the certificate below.

Let $b_r$ and $b_m$ bound changes in the aggregation and merge distance tests respectively. With unchanged radii in fully processed coordinates, both may conservatively be set to $2\varepsilon_x$, by the triangle inequality. If thresholds change, their changes must also be added to the relevant bounds, including the score-window bound.

\begin{proposition}
\label{prop:local}
For a fixed distance-merging configuration and unchanged processed thresholds, the same partition is produced if
\[
 g_s>2b_s,\quad g_o>b_s,\quad g_w>2b_s,
 \quad g_r>b_r,\quad g_m>b_m,
\]
provided that any later nearest-eligible-group choices and tie decisions in small-component handling are also unchanged. The input identities, feature dimension, minimum-size rule and implementation branch are fixed. The claim is in exact arithmetic, or with numerical error included in the bounds.
\end{proposition}
\begin{proof}
The first condition preserves the sorted order: two scores can approach each other by at most $2b_s$. The second preserves the orientation anchor's sign, so the package's orientation agrees with the sign alignment used in the score bound. The third preserves the relevant search windows. Proceeding through the original traversal, the next unassigned starting-point identity is therefore the same; the fourth condition preserves its accept/reject decisions. Induction gives the same aggregation memberships and representative identities, although their coordinates may move. The third and fifth conditions then preserve the distance-merge links, hence their connected components and sizes. Minimum-size eligibility is consequently unchanged. The additional assumption on nearest-group choices and ties preserves any subsequent reallocation, giving the same final partition up to label names.
\end{proof}

Each inequality says that a decision's distance from its boundary must exceed the largest allowed movement. These requirements are deliberately conservative: failure of one does not imply a changed partition. The certificate is local, conditional and sufficient, not necessary. It is not a general proof of CLASSIX stability, and the recorded experiments do not establish that all these margins hold. Density merging additionally depends on overlap counts and volumes, which these conditions do not control. In particular, the selected rich CLASSIX fit uses density merging and is not certified by Proposition~\ref{prop:local}. Its empirical perturbation results in Section~\ref{sec:perturbation_results} are complementary observations, not an application of this proposition.

\subsection{Benchmark and robustness designs}

Nine labelled datasets provide an algorithm-level benchmark for CLASSIX, K Means, DBSCAN and Ward \citep{ester1996dbscan,ward1963hierarchical}. Candidate selection uses only internal indices and deployment guards. Reference labels are reserved for ARI and NMI evaluation. K Means results average five fixed seeds before parameter selection. Ward, DBSCAN and CLASSIX store one structural row per candidate, while repeated fits audit label identity and runtime on the recorded environment. This empirical identity check is not a global determinism theorem.

The rich representation undergoes twelve leave-one-feature-out fits. Principal component analysis (PCA) rotates the standardised features into orthogonal directions ordered by explained variance. Here it retains at least 90 per cent of standardised variance and is compared with the named space on structure and stability, never on rule quality \citep{pearson1901pca,hotelling1933pca}. The CLASSIX perturbation check adds independent zero-mean Gaussian noise to each standardised feature entry, then refits with the selected parameters and compares labels with the unperturbed fit. It uses noise standard deviations 0, 0.000001, 0.0001, 0.001 and 0.01 with ten seeds at every nonzero level. Session counts are rebuilt with gaps of 15, 30 and 60 minutes. A negative control appends independently permuted rich-only dimensions to unchanged compact features. These diagnostics retain the primary selected configurations unless a section explicitly reports complete feature reselection.
